\chapter{Justificación}

El proyecto aborda un problema de investigación en visión por computador médica: aprender representaciones para clasificar lesiones cutáneas cuando las etiquetas clínicas son limitadas. El aprendizaje autosupervisado permite usar imágenes no etiquetadas durante la fase de representación y evaluar después cómo esas representaciones sirven para la tarea de clasificación \parencite{krishnan2022selfsupervisedmedicine}.

La contribución metodológica consiste en documentar y evaluar un clasificador SSL, compararlo con controles explícitos y comunicar sus límites. Los resultados se interpretarán en relación con el protocolo y el conjunto de datos usados; no se presentarán como diagnóstico, eficacia clínica ni mejora asistencial. La evidencia local de carcinoma basocelular ofrece contexto regional, pero no demuestra que HAM10000 represente la epidemiología completa de Santander \parencite{uribe2018bccbucaramanga,tschandl2018ham10000}.

El proyecto también permite estudiar, sin asumir de antemano que aporten al clasificador, módulos complementarios de segmentación, manejo de artefactos y estimación de tono. El flujo modular exige validarlos de forma independiente y medir su aporte incremental antes de integrarlos. Esta separación permite conservar hallazgos negativos o no concluyentes y evita atribuir al SSL mejoras que podrían deberse a cambios en los datos, la partición o el preprocesamiento.

La viabilidad se apoya en conjuntos de datos públicos de investigación, código reproducible y recursos institucionales de cómputo. El trabajo se delimita al experimento computacional y a un prototipo de laboratorio; no incluye recolección propia de imágenes, pruebas con pacientes ni despliegue hospitalario. La representatividad de los conjuntos, la comparabilidad con líneas base supervisadas y la incertidumbre estadística se tratarán como limitaciones, no como beneficios ya demostrados.




